% calcu5.tex — Part I: Prologue (source: calcu5.pdf).
%
% Contains two chapters: 1. Basic Properties of Numbers and 2. Numbers of
% Various Sorts. The source opens with a part divider, then an epigraph, then
% the chapters. Every equation was re-typed; the OCR text layer mangles
% subscripts and symbols (e.g. "a\ + • • • + an" -> $a_1 + \cdots + a_n$).
%
% Items flagged "% [reconstructed]" below are formulas the OCR dropped (they
% sit in the scan as images); they were restored from the standard Spivak
% text and must be checked against the source scan by a vision-capable agent.

\partdivider{I}{Prologue}

\epigraph{%
  To be conscious that you are ignorant is a great step to knowledge.%
}{Benjamin Disraeli}

\spivakchapter{1}{Basic Properties of Numbers}

The title of this chapter expresses in a few words the mathematical knowledge
required to read this book. In fact, this short chapter is simply an
explanation of what is meant by the ``basic properties of numbers,'' all of
which---addition and multiplication, subtraction and division, solutions of
equations and inequalities, factoring and other algebraic manipulations---are
already familiar to us. Nevertheless, this chapter is not a review. Despite
the familiarity of the subject, the survey we are about to undertake will
probably seem quite novel; it does not aim to present an extended review of
old material, but to condense this knowledge into a few simple and obvious
properties of numbers. Some may even seem too obvious to mention, but a
surprising number of diverse and important facts turn out to be consequences
of the ones we shall emphasize.

Of the twelve properties which we shall study in this chapter, the first nine
are concerned with the fundamental operations of addition and multiplication.
For the moment we consider only addition: this operation is performed on a
pair of numbers---the sum $a+b$ exists for any two given numbers $a$ and $b$
(which may possibly be the same number, of course). It might seem reasonable
to regard addition as an operation which can be performed on several numbers
at once, and consider the sum $a_1 + \cdots + a_n$ of $n$ numbers
$a_1, \dots, a_n$ as a basic concept. It is more convenient, however, to
consider addition of pairs of numbers only, and to define other sums in terms
of sums of this type. For the sum of three numbers $a$, $b$, and $c$, this
may be done in two different ways. One can first add $b$ and $c$, obtaining
$b+c$, and then add $a$ to this number, obtaining $a+(b+c)$; or one can first
add $a$ and $b$, and then add the sum $a+b$ to $c$, obtaining $(a+b)+c$. Of
course, the two compound sums obtained are equal, and this fact is the very
first property we shall list:

\begin{quote}
\prop{1}\ \ If $a$, $b$, and $c$ are any numbers, then
\[ a + (b + c) = (a + b) + c. \]
\end{quote}

The statement of this property clearly renders a separate concept of the sum
of three numbers superfluous; we simply agree that $a+b+c$ denotes the number
$a+(b+c)=(a+b)+c$. Addition of four numbers requires similar, though slightly
more involved, considerations. The symbol $a+b+c+d$ is defined to mean

\begin{align*}
&(1) & ((a+b)+c)+d,\\
\text{or } &(2) & (a+(b+c))+d,\\
\text{or } &(3) & a+((b+c)+d),\\
\text{or } &(4) & a+(b+(c+d)),\\
\text{or } &(5) & (a+b)+(c+d).
\end{align*}

This definition is unambiguous since these numbers are all equal.
Fortunately, this fact need not be listed separately, since it follows from
the property P1 already listed. For example, we know from P1 that
\[ (a+b)+c = a+(b+c), \]
and it follows immediately that (1) and (2) are equal. The equality of (2)
and (3) is a direct consequence of P1, although this may not be apparent at
first sight (one must let $b+c$ play the role of $b$ in P1, and $d$ the role
of $c$). The equalities $(3)=(4)=(5)$ are also simple to prove.

It is probably obvious that an appeal to P1 will also suffice to prove the
equality of the 14 possible ways of summing five numbers, but it may not be so
clear how we can reasonably arrange a proof that this is so without actually
listing these 14 sums. Such a procedure is feasible, but would soon cease to
be if we considered collections of six, seven, or more numbers; it would be
totally inadequate to prove the equality of all possible sums of an arbitrary
finite collection of numbers $a_1, \dots, a_n$. This fact may be taken for
granted, but for those who would like to worry about the proof (and it is
worth worrying about once) a reasonable approach is outlined in Problem 24.
Henceforth, we shall usually make a tacit appeal to the results of this
problem and write sums $a_1 + \cdots + a_n$ with a blithe disregard for the
arrangement of parentheses.

The number $0$ has one property so important that we list it next:

\begin{quote}
\prop{2}\ \ If $a$ is any number, then
\[ a + 0 = 0 + a = a. \]
\end{quote}

An important role is also played by $0$ in the third property of our list:

\begin{quote}
\prop{3}\ \ For every number $a$, there is a number $-a$ such that
\[ a + (-a) = (-a) + a = 0. \]
\end{quote}

Property P2 ought to represent a distinguishing characteristic of the number
$0$, and it is comforting to note that we are already in a position to prove
this. Indeed, if a number $x$ satisfies
\[ a + x = a \]
for any one number $a$, then $x=0$ (and consequently this equation also holds
for all numbers $a$). The proof of this assertion involves nothing more than
subtracting $a$ from both sides of the equation, in other words, adding $-a$
to both sides; as the following detailed proof shows, all three properties
P1--P3 must be used to justify this operation.
\begin{align*}
&\text{If} & a + x &= a,\\
&\text{then} & (-a) + (a+x) &= (-a) + a;\\
&\text{hence} & ((-a)+a)+x &= 0;\\
&\text{hence} & 0 + x &= 0;\\
&\text{hence} & x &= 0.
\end{align*}

As we have just hinted, it is convenient to regard subtraction as an
operation derived from addition: we consider $a-b$ to be an abbreviation for
$a+(-b)$. It is then possible to find the solution of certain simple
equations by a series of steps (each justified by P1, P2, or P3) similar to
the ones just presented for the equation $a+x=a$. For example:
\begin{align*}
&\text{If} & x + 3 &= 5,\\
&\text{then} & (x+3)+(-3) &= 5 + (-3);\\
&\text{hence} & x + (3+(-3)) &= 5 - 3 = 2;\\
&\text{hence} & x + 0 &= 2;\\
&\text{hence} & x &= 2.
\end{align*}
Naturally, such elaborate solutions are of interest only until you become
convinced that they can always be supplied. In practice, it is usually just a
waste of time to solve an equation by indicating so explicitly the reliance
on properties P1, P2, and P3 (or any of the further properties we shall
list).

Only one other property of addition remains to be listed. When considering
the sums of three numbers $a$, $b$, and $c$, only two sums were mentioned:
$(a+b)+c$ and $a+(b+c)$. Actually, several other arrangements are obtained if
the order of $a$, $b$, and $c$ is changed. That these sums are all equal
depends on

\begin{quote}
\prop{4}\ \ If $a$ and $b$ are any numbers, then
\[ a + b = b + a. \]
\end{quote}

The statement of P4 is meant to emphasize that although the operation of
addition of pairs of numbers might conceivably depend on the order of the two
numbers, in fact it does not. It is helpful to remember that not all
operations are so well behaved. For example, subtraction does not have this
property: usually $a-b \neq b-a$. In passing we might ask just when $a-b$ does
equal $b-a$, and it is amusing to discover how powerless we are if we rely
only on properties P1--P4 to justify our manipulations. Algebra of the most
elementary variety shows that $a-b=b-a$ only when $a=b$. Nevertheless, it is
impossible to derive this fact from properties P1--P4; it is instructive to
examine the elementary algebra carefully and determine which step(s) cannot
be justified by P1--P4. We will indeed be able to justify all steps in detail
when a few more properties are listed. Oddly enough, however, the crucial
property involves multiplication.

The basic properties of multiplication are fortunately so similar to those
for addition that little comment will be needed; both the meaning and the
consequences should be clear. (As in elementary algebra, the product of $a$
and $b$ will be denoted by $a \cdot b$, or simply $ab$.)

\begin{quote}
\prop{5}\ \ If $a$, $b$, and $c$ are any numbers, then
\[ a \cdot (b \cdot c) = (a \cdot b) \cdot c. \]
\end{quote}

\begin{quote}
\prop{6}\ \ If $a$ is any number, then
\[ a \cdot 1 = 1 \cdot a = a. \]
\end{quote}

Moreover, $1 \neq 0$.

(The assertion that $1 \neq 0$ may seem a strange fact to list, but we have
to list it, because there is no way it could possibly be proved on the basis
of the other properties listed---these properties would all hold if there
were only one number, namely, $0$.)

\begin{quote}
\prop{7}\ \ For every number $a \neq 0$, there is a number $a^{-1}$ such that
\[ a \cdot a^{-1} = a^{-1} \cdot a = 1. \]
\end{quote}

\begin{quote}
\prop{8}\ \ If $a$ and $b$ are any numbers, then
\[ a \cdot b = b \cdot a. \]
\end{quote}

One detail which deserves emphasis is the appearance of the condition
$a \neq 0$ in P7. This condition is quite necessary; since $0 \cdot b = 0$
for all numbers $b$, there is no number $0^{-1}$ satisfying
$0 \cdot 0^{-1} = 1$. This restriction has an important consequence for
division. Just as subtraction was defined in terms of addition, so division
is defined in terms of multiplication: The symbol $a/b$ means $a \cdot b^{-1}$.
Since $0^{-1}$ is meaningless, $a/0$ is also meaningless---division by $0$ is
always undefined.

Property P7 has two important consequences. If $a \cdot b = a \cdot c$, it
does not necessarily follow that $b = c$; for if $a = 0$, then both
$a \cdot b$ and $a \cdot c$ are $0$, no matter what $b$ and $c$ are. However,
if $a \neq 0$, then $b = c$; this can be deduced from P7 as follows:
\begin{align*}
&\text{If} & a \cdot b &= a \cdot c \text{ and } a \neq 0,\\
&\text{then} & a^{-1} \cdot (a \cdot b) &= a^{-1} \cdot (a \cdot c);\\
&\text{hence} & (a^{-1} \cdot a) \cdot b &= (a^{-1} \cdot a) \cdot c;\\
&\text{hence} & 1 \cdot b &= 1 \cdot c;\\
&\text{hence} & b &= c.
\end{align*}
It is also a consequence of P7 that if $a \cdot b = 0$, then either $a = 0$ or
$b = 0$. In fact,
\begin{align*}
&\text{if} & a \cdot b &= 0 \text{ and } a \neq 0,\\
&\text{then} & a^{-1} \cdot (a \cdot b) &= 0;\\
&\text{hence} & (a^{-1} \cdot a) \cdot b &= 0;\\
&\text{hence} & 1 \cdot b &= 0;\\
&\text{hence} & b &= 0.
\end{align*}
(It may happen that $a = 0$ and $b = 0$ are both true; this possibility is
not excluded when we say ``either $a = 0$ or $b = 0$''; in mathematics ``or''
is always used in the sense of ``one or the other, or both.'')

This latter consequence of P7 is constantly used in the solution of
equations. Suppose, for example, that a number $x$ is known to satisfy
\[ (x-1)(x-2) = 0. \]
Then it follows that either $x-1=0$ or $x-2=0$; hence $x=1$ or $x=2$.

On the basis of the eight properties listed so far it is still possible to
prove very little. Listing the next property, which combines the operations
of addition and multiplication, will alter this situation drastically.

\begin{quote}
\prop{9}\ \ If $a$, $b$, and $c$ are any numbers, then
\[ a \cdot (b + c) = a \cdot b + a \cdot c. \]
\end{quote}

(Notice that the equation $(b+c) \cdot a = b \cdot a + c \cdot a$ is also
true, by P8.)

As an example of the usefulness of P9 we will now determine just when
$a-b = b-a$:
\begin{align*}
&\text{If} & a - b &= b - a,\\
&\text{then} & (a-b)+b &= (b-a)+b = b+(b-a);\\
&\text{hence} & a &= b + b - a;\\
&\text{hence} & a + a &= (b+b-a)+a = b+b.\\
&\text{Consequently} & a \cdot (1+1) &= b \cdot (1+1),\\
&\text{and therefore} & a &= b.
\end{align*}

A second use of P9 is the justification of the assertion $a \cdot 0 = 0$
which we have already made, and even used in a proof on page 6 (can you find
where?). This fact was not listed as one of the basic properties, even though
no proof was offered when it was first mentioned. With P1--P8 alone a proof
was not possible, since the number $0$ appears only in P2 and P3, which
concern addition, while the assertion in question involves multiplication.
With P9 the proof is simple, though perhaps not obvious: We have
\begin{align*}
a \cdot 0 + a \cdot 0 &= a \cdot (0 + 0)\\
&= a \cdot 0;
\end{align*}
as we have already noted, this immediately implies (by adding $-(a \cdot 0)$
to both sides) that $a \cdot 0 = 0$.

A series of further consequences of P9 may help explain the somewhat
mysterious rule that the product of two negative numbers is positive. To
begin with, we will establish the more easily acceptable assertion that
$(-a) \cdot b = -(a \cdot b)$. To prove this, note that
\begin{align*}
(-a) \cdot b + a \cdot b &= [(-a)+a] \cdot b\\
&= 0 \cdot b\\
&= 0.
\end{align*}
It follows immediately (by adding $-(a \cdot b)$ to both sides) that
$(-a) \cdot b = -(a \cdot b)$. Now note that
\begin{align*}
(-a) \cdot (-b) + [-(a \cdot b)] &= (-a) \cdot (-b) + (-a) \cdot b\\
&= (-a) \cdot [(-b) + b]\\
&= (-a) \cdot 0\\
&= 0.
\end{align*}
Consequently, adding $(a \cdot b)$ to both sides, we obtain
\[ (-a) \cdot (-b) = a \cdot b. \]
The fact that the product of two negative numbers is positive is thus a
consequence of P1--P9. In other words, \emph{if we want P1 to P9 to be true,
the rule for the product of two negative numbers is forced upon us.}

The various consequences of P9 examined so far, although interesting and
important, do not really indicate the significance of P9; after all, we could
have listed each of these properties separately. Actually, P9 is the
justification for almost all algebraic manipulations. For example, although
we have shown how to solve the equation
\[ (x-1)(x-2) = 0, \]
we can hardly expect to be presented with an equation in this form. We are
more likely to be confronted with the equation
\[ x^2 - 3x + 2 = 0. \]
The ``factorization'' $x^2 - 3x + 2 = (x-1)(x-2)$ is really a triple use of P9:
\begin{align*}
(x-1) \cdot (x-2) &= x \cdot (x-2) + (-1) \cdot (x-2)\\
&= x \cdot x + x \cdot (-2) + (-1) \cdot x + (-1) \cdot (-2)\\
&= x^2 + x[(-2)+(-1)] + 2\\
&= x^2 - 3x + 2.
\end{align*}

A final illustration of the importance of P9 is the fact that this property
is actually used every time one multiplies arabic numerals. For example, the
calculation
\begin{center}
\begin{tabular}{r}
$13$\\
$\times\ 24$\\ \hline
$52$\\
$26$\\ \hline
$312$
\end{tabular}
\end{center}
is a concise arrangement for the following equations:
\begin{align*}
13 \cdot 24 &= 13 \cdot (2 \cdot 10 + 4)\\
&= 13 \cdot 2 \cdot 10 + 13 \cdot 4\\
&= 26 \cdot 10 + 52.
\end{align*}
(Note that moving 26 to the left in the above calculation is the same as
writing $26 \cdot 10$.) The multiplication $13 \cdot 4 = 52$ uses P9 also:
\begin{align*}
13 \cdot 4 &= (1 \cdot 10 + 3) \cdot 4\\
&= 1 \cdot 10 \cdot 4 + 3 \cdot 4\\
&= 4 \cdot 10 + 12\\
&= 4 \cdot 10 + 1 \cdot 10 + 2\\
&= (4+1) \cdot 10 + 2\\
&= 5 \cdot 10 + 2\\
&= 52.
\end{align*}

The properties P1--P9 have descriptive names which are not essential to
remember, but which are often convenient for reference. We will take this
opportunity to list properties P1--P9 together and indicate the names by
which they are commonly designated.

\begin{center}
\renewcommand{\arraystretch}{1.4}
\begin{tabular}{@{}p{0.58\textwidth}p{0.38\textwidth}@{}}
\textbf{(P1)} \textit{(Associative law for addition)} & $a+(b+c)=(a+b)+c$.\\[0.4em]
\textbf{(P2)} \textit{(Existence of an additive identity)} & $a+0=0+a=a$.\\[0.4em]
\textbf{(P3)} \textit{(Existence of additive inverses)} & $a+(-a)=(-a)+a=0$.\\[0.4em]
\textbf{(P4)} \textit{(Commutative law for addition)} & $a+b=b+a$.\\[0.4em]
\textbf{(P5)} \textit{(Associative law for multiplication)} & $a\cdot(b\cdot c)=(a\cdot b)\cdot c$.\\[0.4em]
\textbf{(P6)} \textit{(Existence of a multiplicative identity)} & $a\cdot 1=1\cdot a=a$;\ $1\neq 0$.\\[0.4em]
\textbf{(P7)} \textit{(Existence of multiplicative inverses)} & $a\cdot a^{-1}=a^{-1}\cdot a=1$, for $a\neq 0$.\\[0.4em]
\textbf{(P8)} \textit{(Commutative law for multiplication)} & $a\cdot b=b\cdot a$.\\[0.4em]
\textbf{(P9)} \textit{(Distributive law)} & $a\cdot(b+c)=a\cdot b+a\cdot c$.
\end{tabular}
\end{center}

The three basic properties of numbers which remain to be listed are concerned
with inequalities. Although inequalities occur rarely in elementary
mathematics, they play a prominent role in calculus. The two notions of
inequality, $a<b$ ($a$ is less than $b$) and $a>b$ ($a$ is greater than $b$),
are intimately related: $a<b$ means the same as $b>a$ (thus $1<3$ and $3>1$
are merely two ways of writing the same assertion). The numbers $a$
satisfying $a>0$ are called \emph{positive}, while those numbers $a$
satisfying $a<0$ are called \emph{negative}. While positivity can thus be
defined in terms of $<$, it is possible to reverse the procedure: $a<b$ can
be defined to mean that $b-a$ is positive. In fact, it is convenient to
consider the collection of all positive numbers, denoted by $P$, as the basic
concept, and state all properties in terms of $P$:

\begin{quote}
\prop{10}\ \ \textit{(Trichotomy law)} For every number $a$, one and only
one of the following holds:
\begin{romans}
\item $a = 0$,
\item $a$ is in the collection $P$,
\item $-a$ is in the collection $P$.
\end{romans}
\end{quote}

\begin{quote}
\prop{11}\ \ \textit{(Closure under addition)} If $a$ and $b$ are in $P$,
then $a+b$ is in $P$.
\end{quote}

\begin{quote}
\prop{12}\ \ \textit{(Closure under multiplication)} If $a$ and $b$ are in
$P$, then $a \cdot b$ is in $P$.
\end{quote}

These three properties should be complemented with the following definitions:
\begin{align*}
a > b &\quad\text{if}\quad a - b \text{ is in } P;\\
a < b &\quad\text{if}\quad b > a;\\
a \geq b &\quad\text{if}\quad a > b \text{ or } a = b;\\
a \leq b &\quad\text{if}\quad a < b \text{ or } a = b.
\end{align*}
Note, in particular, that $a>0$ if and only if $a$ is in $P$.

All the familiar facts about inequalities, however elementary they may seem,
are consequences of P10--P12. For example, if $a$ and $b$ are any two
numbers, then precisely one of the following holds:
\begin{romans}
\item $a - b = 0$,
\item $a - b$ is in the collection $P$,
\item $-(a-b) = b - a$ is in the collection $P$.
\end{romans}
Using the definitions just made, it follows that precisely one of the
following holds:
\begin{romans}
\item $a = b$,
\item $a > b$,
\item $b > a$.
\end{romans}

A slightly more interesting fact results from the following manipulations. If
$a<b$, so that $b-a$ is in $P$, then surely $(b+c)-(a+c)$ is in $P$; thus, if
$a<b$, then $a+c<b+c$. Similarly, suppose $a<b$ and $b<c$. Then
\begin{align*}
&b - a \text{ is in } P,\\
&\text{and}\quad c - b \text{ is in } P,\\
&\text{so}\quad c - a = (c-b)+(b-a) \text{ is in } P.
\end{align*}
This shows that if $a<b$ and $b<c$, then $a<c$. (The two inequalities $a<b$
and $b<c$ are usually written in the abbreviated form $a<b<c$, which has the
third inequality $a<c$ almost built in.)

The following assertion is somewhat less obvious: If $a<0$ and $b<0$, then
$ab>0$. The only difficulty presented by the proof is the unraveling of
definitions. The symbol $a<0$ means, by definition, $0>a$, which means
$0-a=-a$ is in $P$. Similarly $-b$ is in $P$, and consequently, by P12,
$(-a)(-b)=ab$ is in $P$. Thus $ab>0$.

The fact that $ab>0$ if $a>0$, $b>0$ and also if $a<0$, $b<0$ has one special
consequence: $a^2>0$ if $a\neq 0$. Thus squares of nonzero numbers are always
positive, and in particular we have proved a result which might have seemed
sufficiently elementary to be included in our list of properties: $1>0$ (since
$1=1^2$).

There is one slightly perplexing feature of the symbols $\geq$ and $\leq$.
The statements
\[ 1+1 \leq 3 \qquad 1+1 \leq 2 \]
are both true, even though we know that $\leq$ could be replaced by $<$ in
the first, and by $=$ in the second. This sort of thing is bound to occur
when $\leq$ is used with specific numbers; the usefulness of the symbol is
revealed by a statement like Theorem 1 where equality holds for some values
of $a$ and $b$, while inequality holds for other values.

The fact that $-a>0$ if $a<0$ is the basis of a concept which will play an
extremely important role in this book. For any number $a$, we define the
absolute value $|a|$ of $a$ as follows:
\[
|a| =
\begin{cases}
a, & a \geq 0\\
-a, & a \leq 0.
\end{cases}
\]
Note that $|a|$ is always positive, except when $a=0$. For example, we have
$|-3|=3$, $|7|=7$, $|1+\sqrt{2}-\sqrt{3}| = 1+\sqrt{2}-\sqrt{3}$, and
$|1+\sqrt{2}-\sqrt{10}| = \sqrt{10}-\sqrt{2}-1$. In general, the most
straightforward approach to any problem involving absolute values requires
treating several cases separately, since absolute values are defined by cases
to begin with. This approach may be used to prove the following very
important fact about absolute values.

\begin{theorem}
For all numbers $a$ and $b$, we have
\[ |a+b| \leq |a| + |b|. \]
\end{theorem}

\begin{proof}
We will consider 4 cases:
\begin{romans}
\item $a \geq 0,\ b \geq 0$;
\item $a \geq 0,\ b \leq 0$;
\item $a \leq 0,\ b \geq 0$;
\item $a \leq 0,\ b \leq 0$.
\end{romans}
In case (1) we also have $a+b \geq 0$, and the theorem is obvious; in fact,
\[ |a+b| = a+b = |a|+|b|, \]
so that in this case equality holds.

In case (4) we have $a+b \leq 0$, and again equality holds:
\[ |a+b| = -(a+b) = -a+(-b) = |a|+|b|. \]

In case (2), when $a \geq 0$ and $b \leq 0$, we must prove that
\[ |a+b| \leq a - b. \]
This case may therefore be divided into two subcases. If $a+b \geq 0$, then we
must prove that
\[ a+b \leq a-b, \]
i.e., $b \leq -b$, which is certainly true since $b \leq 0$ and hence
$-b \geq 0$. On the other hand, if $a+b \leq 0$, we must prove that
\[ -a-b \leq a-b, \]
i.e., $-a \leq a$, which is certainly true since $a \geq 0$ and hence
$-a \leq 0$.

Finally, note that case (3) may be disposed of with no additional work, by
applying case (2) with $a$ and $b$ interchanged.
\end{proof}

Although this method of treating absolute values (separate consideration of
various cases) is sometimes the only approach available, there are often
simpler methods which may be used. In fact, it is possible to give a much
shorter proof of Theorem 1; this proof is motivated by the observation that
\[ |a| = \sqrt{a^2}. \]
(Here, and throughout the book, $\sqrt{x}$ denotes the positive square root
of $x$; this symbol is defined only when $x \geq 0$.) We may now observe that
\begin{align*}
(|a+b|)^2 &= (a+b)^2 = a^2 + 2ab + b^2\\
&\leq a^2 + 2|a| \cdot |b| + b^2\\
&= |a|^2 + 2|a| \cdot |b| + |b|^2\\
&= (|a|+|b|)^2.
\end{align*}
From this we can conclude that $|a+b| \leq |a|+|b|$ because $x^2 < y^2$
implies $x < y$, provided that $x$ and $y$ are both nonnegative; a proof of
this fact is left to the reader (Problem 5).

One final observation may be made about the theorem we have just proved: a
close examination of either proof offered shows that
\[ |a+b| = |a|+|b| \]
if $a$ and $b$ have the same sign (i.e., are both positive or both negative),
or if one of the two is $0$, while
\[ |a+b| < |a|+|b| \]
if $a$ and $b$ are of opposite signs.

We will conclude this chapter with a subtle point, neglected until now, whose
inclusion is required in a conscientious survey of the properties of numbers.
After stating property P9, we proved that $a-b=b-a$ implies $a=b$. The proof
began by establishing that
\[ a \cdot (1+1) = b \cdot (1+1), \]
from which we concluded that $a=b$. This result is obtained from the equation
$a \cdot (1+1) = b \cdot (1+1)$ by dividing both sides by $1+1$. Division by
$0$ should be avoided scrupulously, and it must therefore be admitted that the
validity of the argument depends on knowing that $1+1 \neq 0$. Problem 25 is
designed to convince you that this fact cannot possibly be proved from
properties P1--P9 alone! Once P10, P11, and P12 are available, however, the
proof is very simple: We have already seen that $1>0$; it follows that
$1+1>0$, and in particular $1+1 \neq 0$.

This last demonstration has perhaps only strengthened your feeling that it is
absurd to bother proving such obvious facts, but an honest assessment of our
present situation will help justify serious consideration of such details. In
this chapter we have assumed that numbers are familiar objects, and that
P1--P12 are merely explicit statements of obvious, well-known properties of
numbers. It would be difficult, however, to justify this assumption. Although
one learns how to ``work with'' numbers in school, just what numbers are,
remains rather vague. A great deal of this book is devoted to elucidating the
concept of numbers, and by the end of the book we will have become quite well
acquainted with them. But it will be necessary to work with numbers
throughout the book. It is therefore reasonable to admit frankly that we do
not yet thoroughly understand numbers; we may still say that, in whatever way
numbers are finally defined, they should certainly have properties P1--P12.

Most of this chapter has been an attempt to present convincing evidence that
P1--P12 are indeed basic properties which we should assume in order to deduce
other familiar properties of numbers. Some of the problems (which indicate
the derivation of other facts about numbers from P1--P12) are offered as
further evidence. It is still a crucial question whether P1--P12 actually
account for all properties of numbers. As a matter of fact, we shall soon see
that they do not. In the next chapter the deficiencies of properties P1--P12
will become quite clear, but the proper means for correcting these
deficiencies is not so easily discovered. The crucial additional basic
property of numbers which we are seeking is profound and subtle, quite unlike
P1--P12. The discovery of this crucial property will require all the work of
Part II of this book. In the remainder of Part I we will begin to see why
some additional property is required; in order to investigate this we will
have to consider a little more carefully what we mean by ``numbers.''

\subsection*{PROBLEMS}

\pnum{1.} Prove the following:
\begin{romans}
\item If $ax = a$ for some number $a \neq 0$, then $x = 1$.
\item $x^2 - y^2 = (x-y)(x+y)$.
\item If $x^2 = y^2$, then $x = y$ or $x = -y$.
\item $x^3 - y^3 = (x-y)(x^2+xy+y^2)$.
\item $x^n - y^n = (x-y)(x^{n-1} + x^{n-2}y + \cdots + xy^{n-2} + y^{n-1})$.
\item $x^3 + y^3 = (x+y)(x^2 - xy + y^2)$. (There is a particularly easy way
  to do this, using (iv), and it will show you how to find a factorization
  for $x^n + y^n$ whenever $n$ is odd.)
\end{romans}

\pnum{2.} What is wrong with the following ``proof''? Let $x=y$. Then
\begin{align*}
x^2 &= xy,\\
x^2 - y^2 &= xy - y^2,\\
(x+y)(x-y) &= y(x-y),\\
x+y &= y,\\
2y &= y,\\
2 &= 1.
\end{align*}

\pnum{3.} Prove the following:
\begin{romans}
\item $\dfrac{a}{b} = \dfrac{ac}{bc}$, if $b, c \neq 0$.
\item $\dfrac{a}{b} + \dfrac{c}{d} = \dfrac{ad + bc}{bd}$, if $b, d \neq 0$.
\item $(ab)^{-1} = a^{-1}b^{-1}$, if $a, b \neq 0$. (To do this you must
  remember the defining property of $(ab)^{-1}$.)
\item $\dfrac{a}{b} \cdot \dfrac{c}{d} = \dfrac{ac}{bd}$, if $b, d \neq 0$. % [reconstructed]
\item $\dfrac{a/b}{c/d} = \dfrac{ad}{bc}$, if $b, c, d \neq 0$. % [reconstructed]
\item If $b, d \neq 0$, then $\dfrac{a}{b} = \dfrac{c}{d}$ if and only if
  $ad = bc$. Also determine when $\dfrac{a}{b} = \dfrac{b}{a}$.
\end{romans}

\pnum{4.} Find all numbers $x$ for which
\begin{romans}
\item $4 - x < 3 - 2x$.
\item $5 - x^2 < 8$.
\item $5 - x^2 < -2$.
\item $(x-1)(x-3) > 0$. (When is a product of two numbers positive?)
\item $x^2 - 2x + 2 > 0$.
\item $x^2 + x + 1 > 2$.
\item $x^2 - x + 10 > 16$.
\item $x^2 + x + 1 > 0$.
\item $(x-\pi)(x+5)(x-3) > 0$.
\item $(x-\sqrt{2})(x-\sqrt{2}) > 0$.
\item $2^x < 8$.
\item $x + 3^x < 4$.
\item $\dfrac{1}{x} + \dfrac{1}{1-x} > 0$.
\item $\dfrac{x-1}{x+1} > 0$. % [reconstructed]
\end{romans}

\pnum{5.} Prove the following:
\begin{romans}
\item If $a < b$ and $c < d$, then $a + c < b + d$.
\item If $a < b$, then $-b < -a$.
\item If $a < b$ and $c > d$, then $a - c < b - d$.
\item If $a < b$ and $c > 0$, then $ac < bc$.
\item If $a < b$ and $c < 0$, then $ac > bc$.
\item If $a > 1$, then $a^2 > a$.
\item If $0 < a < 1$, then $a^2 < a$.
\item If $0 \leq a < b$ and $0 \leq c < d$, then $ac < bd$.
\item If $0 \leq a < b$, then $a^2 < b^2$. (Use (viii).)
\item If $a, b \geq 0$ and $a^2 < b^2$, then $a < b$. (Use (ix), backwards.)
\end{romans}

\pnum{6.}
\begin{letters}
\item Prove that if $0 \leq x < y$, then $x^n < y^n$, $n = 1, 2, 3, \dots$.
\item Prove that if $x < y$ and $n$ is odd, then $x^n < y^n$.
\item Prove that if $x^n = y^n$ and $n$ is odd, then $x = y$.
\item Prove that if $x^n = y^n$ and $n$ is even, then $x = y$ or $x = -y$.
\end{letters}

\pnum{7.} Prove that if $0 < a < b$, then
\[ a < \sqrt{ab} < \frac{a+b}{2} < b. \]
Notice that the inequality $\sqrt{ab} \leq (a+b)/2$ holds for all $a, b \geq 0$.
A generalization of this fact occurs in Problem 2-22.

\pnum{*8.} Although the basic properties of inequalities were stated in terms
of the collection $P$ of all positive numbers, and $<$ was defined in terms of
$P$, this procedure can be reversed. Suppose that P10--P12 are replaced by
\begin{quote}
\textbf{(P$'$10)}\ \ For any numbers $a$ and $b$ one, and only one, of the
following holds:
\begin{romans}
\item $a = b$,
\item $a < b$,
\item $b < a$.
\end{romans}

\textbf{(P$'$11)}\ \ For any numbers $a$, $b$, and $c$, if $a < b$ and
$b < c$, then $a < c$.

\textbf{(P$'$12)}\ \ For any numbers $a$, $b$, and $c$, if $a < b$, then
$a + c < b + c$.

\textbf{(P$'$13)}\ \ For any numbers $a$, $b$, and $c$, if $a < b$ and
$0 < c$, then $ac < bc$.
\end{quote}
Show that P10--P12 can then be deduced as theorems.

\pnum{9.} Express each of the following with at least one less pair of
absolute value signs.
\begin{romans}
\item $|\sqrt{2} + \sqrt{3} - \sqrt{5} + \sqrt{7}|$.
\item $|(|a+b| - |a| - |b|)|$.
\item $|(|a+b| + |c| - |a+b+c|)|$.
\item $|x^2 - 2xy + y^2|$.
\item $|(|\sqrt{2} + \sqrt{3}| - |\sqrt{5} - \sqrt{7}|)|$.
\end{romans}

\pnum{10.} Express each of the following without absolute value signs,
treating various cases separately when necessary.
\begin{romans}
\item $|a+b| - |b|$.
\item $|(|x|-1)|$.
\item $|x| - |x^2|$.
\item $a - |(a-|a|)|$.
\end{romans}

\pnum{11.} Find all numbers $x$ for which
\begin{romans}
\item $|x-3| = 8$.
\item $|x-3| < 8$.
\item $|x+4| < 2$.
\item $|x-1| + |x-2| > 1$.
\item $|x-1| + |x+1| < 2$.
\item $|x-1| + |x+1| < 1$.
\item $|x-1| \cdot |x+1| = 0$.
\item $|x-1| \cdot |x+2| = 3$.
\end{romans}

\pnum{12.} Prove the following:
\begin{romans}
\item $|xy| = |x| \cdot |y|$.
\item $\dfrac{1}{|x|} = \left|\dfrac{1}{x}\right|$, if $x \neq 0$. (The best
  way to do this is to remember what $|x|^{-1}$ is.)
\item $\dfrac{|x|}{|y|} = \left|\dfrac{x}{y}\right|$, if $y \neq 0$.
\item $|x-y| \leq |x| + |y|$. (Give a very short proof.)
\item $|x| - |y| \leq |x-y|$. (A very short proof is possible, if you write
  things in the right way.)
\item $|(|x|-|y|)| \leq |x-y|$. (Why does this follow immediately from (v)?)
\item $|x+y+z| \leq |x|+|y|+|z|$. Indicate when equality holds, and prove
  your statement.
\end{romans}

\pnum{13.} The maximum of two numbers $x$ and $y$ is denoted by $\max(x,y)$.
Thus $\max(-1,3) = \max(3,3) = 3$ and $\max(-1,-4) = \max(-4,-1) = -1$. The
minimum of $x$ and $y$ is denoted by $\min(x,y)$. Prove that
\[
\max(x,y) = \frac{x + y + |y-x|}{2},
\qquad
\min(x,y) = \frac{x + y - |y-x|}{2}.
\]
Derive a formula for $\max(x,y,z)$ and $\min(x,y,z)$, using, for example
\[ \max(x,y,z) = \max(x,\max(y,z)). \]

\pnum{14.}
\begin{letters}
\item Prove that $|a| = |-a|$. (The trick is not to become confused by too
  many cases. First prove the statement for $a \geq 0$. Why is it then obvious
  for $a \leq 0$?)
\item Prove that $-b \leq a \leq b$ if and only if $|a| \leq b$. In
  particular, it follows that $-|a| \leq a \leq |a|$.
\item Use this fact to give a new proof that $|a+b| \leq |a|+|b|$.
\end{letters}

\pnum{*15.} Prove that if $x$ and $y$ are not both $0$, then
\[
x^2 + xy + y^2 > 0,
\qquad
x^4 + x^3y + x^2y^2 + xy^3 + y^4 > 0.
\]
\textit{Hint:} Use Problem 1.

\pnum{*16.}
\begin{letters}
\item Show that
\[
(x+y)^2 = x^2 + y^2 \text{ only when } x = 0 \text{ or } y = 0,
\]
\[
(x+y)^3 = x^3 + y^3 \text{ only when } x = 0 \text{ or } y = 0 \text{ or }
x = -y.
\]
\item Using the fact that
\[ x^2 + 2xy + y^2 = (x+y)^2 \geq 0, \]
show that $4x^2 + 6xy + 4y^2 > 0$ unless $x$ and $y$ are both $0$.
\item Use part (b) to find out when $(x+y)^4 = x^4 + y^4$.
\item Find out when $(x+y)^5 = x^5 + y^5$. \textit{Hint:} From the assumption
  $(x+y)^5 = x^5 + y^5$ you should be able to derive the equation
  $x^2 + 2x^2y + 2xy^2 + y^3 = 0$, if $xy \neq 0$. This implies that
  $(x+y)^3 = x^2y + xy^2 = xy(x+y)$.
\end{letters}
You should now be able to make a good guess as to when $(x+y)^n = x^n + y^n$;
the proof is contained in Problem 11-63.

\pnum{17.}
\begin{letters}
\item Find the smallest possible value of $2x^2 - 3x + 4$. \textit{Hint:}
  ``Complete the square,'' i.e., write $2x^2 - 3x + 4 = 2(x - 3/4)^2 + \text{?}$
\item Find the smallest possible value of $x^2 - 3x + 2y^2 + 4y + 2$.
\item Find the smallest possible value of $x^2 + 4xy + 5y^2 - 4x - 6y + 7$.
\end{letters}

\pnum{18.}
\begin{letters}
\item Suppose that $b^2 - 4c \geq 0$. Show that the numbers
\[ \frac{-b + \sqrt{b^2 - 4c}}{2}, \qquad \frac{-b - \sqrt{b^2 - 4c}}{2} \]
both satisfy the equation $x^2 + bx + c = 0$.
\item Suppose that $b^2 - 4c < 0$. Show that there are no numbers $x$
  satisfying $x^2 + bx + c = 0$; in fact, $x^2 + bx + c > 0$ for all $x$.
  \textit{Hint:} Complete the square.
\item Use this fact to give another proof that if $x$ and $y$ are not both
  $0$, then $x^2 + xy + y^2 > 0$.
\item For which numbers $a$ is it true that $x^2 + axy + y^2 > 0$ whenever
  $x$ and $y$ are not both $0$?
\item Find the smallest possible value of $x^2 + bx + c$ and of
  $ax^2 + bx + c$, for $a > 0$.
\end{letters}

\pnum{19.} The fact that $a^2 \geq 0$ for all numbers $a$, elementary as it
may seem, is nevertheless the fundamental idea upon which most important
inequalities are ultimately based. The great-granddaddy of all inequalities
is the \emph{Schwarz inequality}.
\[ x_1y_1 + x_2y_2 \leq \sqrt{x_1^{\,2} + x_2^{\,2}}\ \sqrt{y_1^{\,2} + y_2^{\,2}}. \]
(A more general form occurs in Problem 2-21.) The three proofs of the Schwarz
inequality outlined below have only one thing in common---their reliance on
the fact that $a^2 \geq 0$ for all $a$.

\begin{letters}
\item Prove that if $x_1 = \lambda y_1$ and $x_2 = \lambda y_2$ for some
  number $\lambda \geq 0$, then equality holds in the Schwarz inequality.
  Prove the same thing if $y_1 = y_2 = 0$. Now suppose that $y_1$ and $y_2$
  are not both $0$, and that there is no number $\lambda$ such that
  $x_1 = \lambda y_1$ and $x_2 = \lambda y_2$. Then
  \begin{align*}
  0 &< (\lambda y_1 - x_1)^2 + (\lambda y_2 - x_2)^2\\
    &= \lambda^2(y_1^{\,2} + y_2^{\,2}) - 2\lambda(x_1y_1 + x_2y_2)
       + (x_1^{\,2} + x_2^{\,2}).
  \end{align*}
  Using Problem 18, complete the proof of the Schwarz inequality.
\item Prove the Schwarz inequality by using $2xy \leq x^2 + y^2$ (how is this
  derived?) with
  \[
  x = \frac{x_i}{\sqrt{x_1^{\,2} + x_2^{\,2}}},
  \qquad
  y = \frac{y_i}{\sqrt{y_1^{\,2} + y_2^{\,2}}},
  \]
  first for $i = 1$ and then for $i = 2$.
\item Prove the Schwarz inequality by first proving that
  \[
  (x_1^{\,2} + x_2^{\,2})(y_1^{\,2} + y_2^{\,2})
  = (x_1y_1 + x_2y_2)^2 + (x_1y_2 - x_2y_1)^2.
  \]
\item Deduce, from each of these three proofs, that equality holds only when
  $y_1 = y_2 = 0$ or when there is a number $\lambda \geq 0$ such that
  $x_1 = \lambda y_1$ and $x_2 = \lambda y_2$.
\end{letters}

In our later work, three facts about inequalities will be crucial. Although
proofs will be supplied at the appropriate point in the text, a personal
assault on these problems is infinitely more enlightening than a perusal of a
completely worked-out proof. The statements of these propositions involve
some weird numbers, but their basic message is very simple: if $x$ is close
enough to $x_0$, and $y$ is close enough to $y_0$, then $x+y$ will be close to
$x_0 + y_0$, and $xy$ will be close to $x_0y_0$, and $1/y$ will be close to
$1/y_0$. The symbol ``$\varepsilon$'' which appears in these propositions is
the fifth letter of the Greek alphabet (``epsilon''), and could just as well
be replaced by a less intimidating Roman letter; however, tradition has made
the use of $\varepsilon$ almost sacrosanct in the contexts to which these
theorems apply.

\pnum{20.} Prove that if
\[
|x - x_0| < \frac{\varepsilon}{2}
\qquad\text{and}\qquad
|y - y_0| < \frac{\varepsilon}{2},
\]
then
\[
|(x+y) - (x_0+y_0)| < \varepsilon,
\qquad
|(x-y) - (x_0-y_0)| < \varepsilon.
\]

\pnum{*21.} Prove that if
\[
|x - x_0| < \min\left(\frac{\varepsilon}{2(|y_0|+1)}, 1\right)
\qquad\text{and}\qquad
|y - y_0| < \frac{\varepsilon}{2(|x_0|+1)}
\]
then $|xy - x_0y_0| < \varepsilon$.
(The notation ``min'' was defined in Problem 13, but the formula provided by
that problem is irrelevant at the moment; the first inequality in the
hypothesis just means that
\[
|x - x_0| < \frac{\varepsilon}{2(|y_0|+1)}
\qquad\text{and}\qquad
|x - x_0| < 1;
\]
at one point in the argument you will need the first inequality, and at
another point you will need the second. One more word of advice: since the
hypotheses only provide information about $x - x_0$ and $y - y_0$, it is
almost a foregone conclusion that the proof will depend upon writing
$xy - x_0y_0$ in a way that involves $x - x_0$ and $y - y_0$.)

\pnum{*22.} Prove that if $y_0 \neq 0$ and
\[
|y - y_0| < \min\left(\frac{|y_0|}{2}, \frac{\varepsilon |y_0|^2}{2}\right)
\]
then $y \neq 0$ and
\[
\left|\frac{1}{y} - \frac{1}{y_0}\right| < \varepsilon.
\]

\pnum{*23.} Replace the question marks in the following statement by
expressions involving $\varepsilon$, $x_0$, and $y_0$ so that the conclusion
will be true:
If $y_0 \neq 0$ and
\[
|y - y_0| < \text{?} \qquad\text{and}\qquad |x - x_0| < \text{?}
\]
then $y \neq 0$ and
\[
\left|\frac{x}{y} - \frac{x_0}{y_0}\right| < \varepsilon.
\]
This problem is trivial in the sense that its solution follows from Problems
21 and 22 with almost no work at all (notice that $x/y = x \cdot 1/y$). The
crucial point is not to become confused; decide which of the two problems
should be used first, and don't panic if your answer looks unlikely.

\pnum{*24.} This problem shows that the actual placement of parentheses in a
sum is irrelevant. The proofs involve ``mathematical induction''; if you are
not familiar with such proofs, but still want to tackle this problem, it can
be saved until after Chapter 2, where proofs by induction are explained.

Let us agree, for definiteness, that $a_1 + \cdots + a_n$ will denote
\[
a_1 + (a_2 + (a_3 + \cdots + (a_{n-2} + (a_{n-1} + a_n))) \cdots).
\]
Thus $a_1 + a_2 + a_3$ denotes $a_1 + (a_2 + a_3)$, and $a_1 + a_2 + a_3 + a_4$
denotes $a_1 + (a_2 + (a_3 + a_4))$, etc.

\begin{letters}
\item Prove that
\[
(a_1 + \cdots + a_k) + a_{k+1} = a_1 + \cdots + a_{k+1}.
\]
\textit{Hint:} Use induction on $k$.
\item Prove that if $n \geq k$, then
\[
(a_1 + \cdots + a_k) + (a_{k+1} + \cdots + a_n) = a_1 + \cdots + a_n.
\]
\textit{Hint:} Use part (a) to give a proof by induction on $k$.
\item Let $s(a_1, \dots, a_k)$ be some sum formed from $a_1, \dots, a_k$.
  Show that
\[ s(a_1, \dots, a_k) = a_1 + \cdots + a_k. \]
\textit{Hint:} There must be two sums $s'(a_1, \dots, a_l)$ and
$s''(a_{l+1}, \dots, a_k)$ such that
\[ s(a_1, \dots, a_k) = s'(a_1, \dots, a_l) + s''(a_{l+1}, \dots, a_k). \]
\end{letters}

\pnum{25.} Suppose that we interpret ``number'' to mean either $0$ or $1$, and
$+$ and $\cdot$ to be the operations defined by the following two tables.
\begin{center}
\begin{tabular}{c|cc}
$+$ & $0$ & $1$\\ \hline
$0$ & $0$ & $1$\\
$1$ & $1$ & $0$
\end{tabular}
\qquad
\begin{tabular}{c|cc}
$\cdot$ & $0$ & $1$\\ \hline
$0$ & $0$ & $0$\\
$1$ & $0$ & $1$
\end{tabular}
\end{center}
Check that properties P1--P9 all hold, even though $1 + 1 = 0$.

\spivakchapter{2}{Numbers of Various Sorts}

In Chapter 1 we used the word ``number'' very loosely, despite our concern
with the basic properties of numbers. It will now be necessary to distinguish
carefully various kinds of numbers.

The simplest numbers are the ``counting numbers''
\[ 1, 2, 3, \dots . \]
The fundamental significance of this collection of numbers is emphasized by
its symbol $\mathbb{N}$ (for \emph{natural} numbers). A brief glance at
P1--P12 will show that our basic properties of ``numbers'' do not apply to
$\mathbb{N}$---for example, P2 and P3 do not make sense for $\mathbb{N}$. From
this point of view the system $\mathbb{N}$ has many deficiencies.
Nevertheless, $\mathbb{N}$ is sufficiently important to deserve several
comments before we consider larger collections of numbers.

The most basic property of $\mathbb{N}$ is the principle of ``mathematical
induction.'' Suppose $P(x)$ means that the property $P$ holds for the number
$x$. Then the principle of mathematical induction states that $P(x)$ is true
for all natural numbers $x$ provided that

\begin{enumerate}[label=(\arabic*),leftmargin=2.6em,itemsep=0pt,topsep=2pt]
\item $P(1)$ is true.
\item Whenever $P(k)$ is true, $P(k+1)$ is true.
\end{enumerate}

Note that condition (2) merely asserts the truth of $P(k+1)$ under the
assumption that $P(k)$ is true; this suffices to ensure the truth of $P(x)$
for all $x$, if condition (1) also holds. In fact, if $P(1)$ is true, then it
follows that $P(2)$ is true (by using (2) in the special case $k = 1$). Now,
since $P(2)$ is true it follows that $P(3)$ is true (using (2) in the special
case $k = 2$). It is clear that each number will eventually be reached by a
series of steps of this sort, so that $P(k)$ is true for all numbers $k$.

A favorite illustration of the reasoning behind mathematical induction
envisions an infinite line of people,
\[ \text{person number 1, person number 2, person number 3, } \dots . \]
If each person has been instructed to tell any secret he hears to the person
behind him (the one with the next largest number) and a secret is told to
person number 1, then clearly every person will eventually learn the secret.
If $P(x)$ is the assertion that person number $x$ will learn the secret, then
the instructions given (to tell all secrets learned to the next person)
assures that condition (2) is true, and telling the secret to person number 1
makes (1) true. The following example is a less facetious use of mathematical
induction. There is a useful and striking formula which expresses the sum of
the first $n$ numbers in a simple way:
\[
1 + \cdots + n = \frac{n(n+1)}{2}.
\]
To prove this formula, note first that it is clearly true for $n = 1$. Now
assume that for some natural number $k$ we have
\[
1 + \cdots + k = \frac{k(k+1)}{2}.
\]
Then
\begin{align*}
1 + \cdots + k + (k+1) &= \frac{k(k+1)}{2} + k + 1\\
&= \frac{k(k+1) + 2k + 2}{2}\\
&= \frac{k^2 + 3k + 2}{2}\\
&= \frac{(k+1)(k+2)}{2},
\end{align*}
so the formula is also true for $k + 1$. By the principle of induction this
proves the formula for all natural numbers $n$. This particular example
illustrates a phenomenon that frequently occurs, especially in connection
with formulas like the one just proved. Although the proof by induction is
often quite straightforward, the method by which the formula was discovered
remains a mystery. Problems 5 and 6 indicate how some formulas of this type
may be derived.

The principle of mathematical induction may be formulated in an equivalent
way without speaking of ``properties'' of a number, a term which is
sufficiently vague to be eschewed in a mathematical discussion. A more precise
formulation states that if $A$ is any collection (or ``set''---a synonymous
mathematical term) of natural numbers and

\begin{enumerate}[label=(\arabic*),leftmargin=2.6em,itemsep=0pt,topsep=2pt]
\item $1$ is in $A$,
\item $k + 1$ is in $A$ whenever $k$ is in $A$,
\end{enumerate}

then $A$ is the set of all natural numbers. It should be clear that this
formulation adequately replaces the less formal one given previously---we just
consider the set $A$ of natural numbers $x$ which satisfy $P(x)$. For example,
suppose $A$ is the set of natural numbers $n$ for which it is true that
\[
1 + \cdots + n = \frac{n(n+1)}{2}.
\]
Our previous proof of this formula showed that $A$ contains 1, and that $k+1$
is in $A$, if $k$ is. It follows that $A$ is the set of all natural numbers,
i.e., that the formula holds for all natural numbers $n$.

There is yet another rigorous formulation of the principle of mathematical
induction, which looks quite different. If $A$ is any collection of natural
numbers, it is tempting to say that $A$ must have a smallest member.
Actually, this statement can fail to be true in a rather subtle way. A
particularly important set of natural numbers is the collection $A$ that
contains no natural numbers at all, the ``empty collection'' or ``null
set,''$^{*}$ denoted by $\emptyset$. The null set $\emptyset$ is a collection
of natural numbers that has no smallest member---in fact, it has no members
at all. This is the only possible exception, however; if $A$ is a nonnull set
of natural numbers, then $A$ has a least member. This ``intuitively obvious''
statement, known as the ``well-ordering principle,'' can be proved from the
principle of induction as follows. Suppose that the set $A$ has no least
member. Let $B$ be the set of natural numbers $n$ such that $1, \dots, n$ are
all not in $A$. Clearly $1$ is in $B$ (because if $1$ were in $A$, then $A$
would have $1$ as smallest member). Moreover, if $1, \dots, k$ are not in $A$,
surely $k+1$ is not in $A$ (otherwise $k+1$ would be the smallest member of
$A$), so $1, \dots, k+1$ are all not in $A$. This shows that if $k$ is in $B$,
then $k+1$ is in $B$. It follows that every number $n$ is in $B$, i.e., the
numbers $1, \dots, n$ are not in $A$ for any natural number $n$. Thus
$A = \emptyset$, which completes the proof.

It is also possible to prove the principle of induction from the well-ordering
principle (Problem 10). Either principle may be considered as a basic
assumption about the natural numbers.

There is still another form of induction which should be mentioned. It
sometimes happens that in order to prove $P(k+1)$ we must assume not only
$P(k)$, but also $P(l)$ for all natural numbers $l < k$. In this case we rely
on the ``principle of complete induction'': If $A$ is a set of natural numbers
and

\begin{enumerate}[label=(\arabic*),leftmargin=2.6em,itemsep=0pt,topsep=2pt]
\item $1$ is in $A$,
\item $k + 1$ is in $A$ if $1, \dots, k$ are in $A$,
\end{enumerate}

then $A$ is the set of all natural numbers.

Although the principle of complete induction may appear much stronger than
the ordinary principle of induction, it is actually a consequence of that
principle. The proof of this fact is left to the reader, with a hint
(Problem 11). Applications will be found in Problems 7, 17, 20 and 22.

Closely related to proofs by induction are ``recursive definitions.'' For
example, the number $n!$ (read ``$n$ factorial'') is defined as the product of
all the natural numbers less than or equal to $n$:
\[ n! = 1 \cdot 2 \cdot \ldots \cdot (n-1) \cdot n. \]
This can be expressed more precisely as follows:
\begin{align*}
&(1) & 1! &= 1\\
&(2) & n! &= n \cdot (n-1)!.
\end{align*}
This form of the definition exhibits the relationship between $n!$ and
$(n-1)!$ in an explicit way that is ideally suited for proofs by induction.
Problem 23 reviews a definition already familiar to you, which may be
expressed more succinctly as a recursive definition; as this problem shows,
the recursive definition is really necessary for a rigorous proof of some of
the basic properties of the definition.

$^{*}$ Although it may not strike you as a collection, in the ordinary sense of
the word, the null set arises quite naturally in many contexts. We frequently
consider the set $A$, consisting of all $x$ satisfying some property $P$;
often we have no guarantee that $P$ is satisfied by any number, so that $A$
might be $\emptyset$---in fact often one proves that $P$ is always false by
showing that $A = \emptyset$.

One definition which may not be familiar involves some convenient notation
which we will constantly be using. Instead of writing
\[ a_1 + \cdots + a_n, \]
we will usually employ the Greek letter $\Sigma$ (capital sigma, for ``sum'')
and write
\[ \sum_{i=1}^{n} a_i. \]
In other words, $\sum_{i=1}^{n} a_i$ denotes the sum of the numbers obtained
by letting $i = 1, 2, \dots, n$. Thus
\[
\sum_{i=1}^{n} a_i
= a_1 + a_2 + \cdots + a_n.
\]
Notice that the letter $i$ really has nothing to do with the number denoted by
$\sum_{i=1}^{n} a_i$, and can be replaced by any convenient symbol (except $n$,
of course!):
\[
\sum_{i=1}^{n} i = \sum_{j=1}^{n} j = \frac{n(n+1)}{2}.
\]
To define $\sum_{i=1}^{n} a_i$ precisely really requires a recursive
definition:
\begin{align*}
&(1) & \sum_{i=1}^{1} a_i &= a_1\\
&(2) & \sum_{i=1}^{n} a_i &= \sum_{i=1}^{n-1} a_i + a_n.
\end{align*}
But only purveyors of mathematical austerity would insist too strongly on
such precision. In practice, all sorts of modifications of this symbolism are
used, and no one ever considers it necessary to add any words of explanation.
The symbol
\[ \sum_{\substack{i=1\\ i \neq 4}}^{n} a_i, \]
for example, is an obvious way of writing
\[ a_1 + a_2 + a_3 + a_5 + \cdots + a_n, \]
or more precisely,
\[ \sum_{i=1}^{3} a_i + \sum_{i=5}^{n} a_i. \]

The deficiencies of the natural numbers which we discovered at the beginning
of this chapter may be partially remedied by extending this system to the set
of integers
\[ \dots, -2, -1, 0, 1, 2, \dots . \]
This set is denoted by $\mathbb{Z}$ (from German ``Zahl,'' number). Of
properties P1--P12, only P7 fails for $\mathbb{Z}$.

A still larger system of numbers is obtained by taking quotients $m/n$ of
integers (with $n \neq 0$). These numbers are called rational numbers, and the
set of all rational numbers is denoted by $\mathbb{Q}$ (for ``quotients''). In
this system of numbers all of P1--P12 are true. It is tempting to conclude
that the ``properties of numbers,'' which we studied in some detail in
Chapter 1, refer to just one set of numbers, namely, $\mathbb{Q}$. There is,
however, a still larger collection of numbers to which properties P1--P12
apply---the set of all real numbers, denoted by $\mathbb{R}$. The real numbers
include not only the rational numbers, but other numbers as well (the
irrational numbers) which can be represented by infinite decimals; $\pi$ and
$\sqrt{2}$ are both examples of irrational numbers. The proof that $\pi$ is
irrational is not easy---we shall devote all of Chapter 16 of Part III to a
proof of this fact. The irrationality of $\sqrt{2}$, on the other hand, is
quite simple, and was known to the Greeks. (Since the Pythagorean theorem
shows that an isosceles right triangle, with sides of length $1$, has a
hypotenuse of length $\sqrt{2}$, it is not surprising that the Greeks should
have investigated this question.) The proof depends on a few observations
about the natural numbers. Every natural number $n$ can be written either in
the form $2k$ for some integer $k$, or else in the form $2k + 1$ for some
integer $k$ (this ``obvious'' fact has a simple proof by induction
(Problem 8)). Those natural numbers of the form $2k$ are called \emph{even};
those of the form $2k+1$ are called \emph{odd}. Note that even numbers have
even squares, and odd numbers have odd squares:
\begin{align*}
(2k)^2 &= 4k^2 = 2 \cdot (2k^2),\\
(2k+1)^2 &= 4k^2 + 4k + 1 = 2 \cdot (2k^2 + 2k) + 1.
\end{align*}
In particular it follows that the converse must also hold: if $n^2$ is even,
then $n$ is even; if $n^2$ is odd, then $n$ is odd. The proof that $\sqrt{2}$
is irrational is now quite simple. Suppose that $\sqrt{2}$ were rational; that
is, suppose there were natural numbers $p$ and $q$ such that
\[ \left(\frac{p}{q}\right)^2 = 2. \]
We can assume that $p$ and $q$ have no common divisor (since all common
divisors could be divided out to begin with). Now we have
\[ p^2 = 2q^2. \]
This shows that $p^2$ is even, and consequently $p$ must be even; that is,
$p = 2k$ for some natural number $k$. Then
\[ p^2 = 4k^2 = 2q^2, \]
so
\[ 2k^2 = q^2. \]
This shows that $q^2$ is even, and consequently that $q$ is even. Thus both
$p$ and $q$ are even, contradicting the fact that $p$ and $q$ have no common
divisor. This contradiction completes the proof.

It is important to understand precisely what this proof shows. We have
demonstrated that there is no rational number $x$ such that $x^2 = 2$. This
assertion is often expressed more briefly by saying that $\sqrt{2}$ is
irrational. Note, however, that the use of the symbol $\sqrt{2}$ implies the
existence of some number (necessarily irrational) whose square is 2. We have
not proved that such a number exists and we can assert confidently that, at
present, a proof is impossible for us. Any proof at this stage would have to
be based on P1--P12 (the only properties of $\mathbb{R}$ we have mentioned);
since P1--P12 are also true for $\mathbb{Q}$ the exact same argument would
show that there is a rational number whose square is 2, and this we know is
false. (Note that the reverse argument will not work---our proof that there is
no rational number whose square is 2 cannot be used to show that there is no
real number whose square is 2, because our proof used not only P1--P12 but
also a special property of $\mathbb{Q}$, the fact that every number in
$\mathbb{Q}$ can be written $p/q$ for integers $p$ and $q$.)

This particular deficiency in our list of properties of the real numbers
could, of course, be corrected by adding a new property which asserts the
existence of square roots of positive numbers. Resorting to such a measure
is, however, neither aesthetically pleasing nor mathematically satisfactory;
we would still not know that every number has an $n$th root if $n$ is odd, and
that every positive number has an $n$th root if $n$ is even. Even if we
assumed this, we could not prove the existence of a number $x$ satisfying
$x^5 + x + 1 = 0$ (even though there does happen to be one), since we do not
know how to write the solution of the equation in terms of $n$th roots (in
fact, it is known that the solution cannot be written in this form). And, of
course, we certainly do not wish to assume that all equations have solutions,
since this is false (no real number $x$ satisfies $x^2 + 1 = 0$, for example).
In fact, this direction of investigation is not a fruitful one. The most
useful hints about the property distinguishing $\mathbb{R}$ from
$\mathbb{Q}$, the most compelling evidence for the necessity of elucidating
this property, do not come from the study of numbers alone. In order to study
the properties of the real numbers in a more profound way, we must study more
than the real numbers. At this point we must begin with the foundations of
calculus, in particular the fundamental concept on which calculus is
based---functions.

\subsection*{PROBLEMS}

\pnum{1.} Prove the following formulas by induction.
\begin{romans}
\item $1^2 + \cdots + n^2 = \dfrac{n(n+1)(2n+1)}{6}$.
\item $1^3 + \cdots + n^3 = (1 + \cdots + n)^2$.
\end{romans}

\pnum{2.} Find a formula for
\begin{romans}
\item $\displaystyle\sum_{i=1}^{n} (2i - 1) = 1 + 3 + 5 + \cdots + (2n - 1)$.
\item $\displaystyle\sum_{i=1}^{n} (2i - 1)^2 = 1^2 + 3^2 + 5^2 + \cdots + (2n - 1)^2$.
\end{romans}
\textit{Hint:} What do these expressions have to do with
$1 + 2 + 3 + \cdots + 2n$ and $1^2 + 2^2 + 3^2 + \cdots + (2n)^2$?

\pnum{3.} If $0 \leq k \leq n$, the ``binomial coefficient'' $\binom{n}{k}$ is
defined by
\[
\binom{n}{k} = \frac{n!}{k!(n-k)!}
= \frac{n(n-1)\cdots(n-k+1)}{k!},
\qquad k \neq 0, n
\]
\[
\binom{n}{0} = \binom{n}{n} = 1,
\]
(a special case of the first formula if we define $0! = 1$), and for $k < 0$
or $k > n$ we just define the binomial coefficient to be $0$.

\begin{letters}
\item Prove that
\[ \binom{n+1}{k} = \binom{n}{k-1} + \binom{n}{k}. \]
(The proof does not require an induction argument.)

This relation gives rise to the following configuration, known as ``Pascal's
triangle''---a number not on one of the sides is the sum of the two numbers
above it; the binomial coefficient $\binom{n}{k}$ is the $(k+1)$st number in
the $(n+1)$st row.
\begin{center}
\begin{tabular}{ccccccccccc}
& & & & & $1$\\
& & & & $1$ & & $1$\\
& & & $1$ & & $2$ & & $1$\\
& & $1$ & & $3$ & & $3$ & & $1$\\
& $1$ & & $4$ & & $6$ & & $4$ & & $1$\\
$1$ & & $5$ & & $10$ & & $10$ & & $5$ & & $1$
\end{tabular}
\end{center}

\item Notice that all the numbers in Pascal's triangle are natural numbers.
  Use part (a) to prove by induction that $\binom{n}{k}$ is always a natural
  number. (Your entire proof by induction will, in a sense, be summed up in a
  glance by Pascal's triangle.)
\item Give another proof that $\binom{n}{k}$ is a natural number by showing
  that $\binom{n}{k}$ is the number of sets of exactly $k$ integers each
  chosen from $1, \dots, n$.
\item Prove the ``binomial theorem'': If $a$ and $b$ are any numbers and $n$
  is a natural number, then
\[ (a+b)^n = a^n + \binom{n}{1} a^{n-1} b + \binom{n}{2} a^{n-2} b^2 + \cdots
   + \binom{n}{n-1} a b^{n-1} + b^n. \]
\item Prove that
\[ \binom{n}{0} + \binom{n}{1} + \cdots + \binom{n}{n} = 2^n. \]
\end{letters}

\pnum{4.}
\begin{letters}
\item Prove that
\[
\sum_{j=0}^{k} \binom{m}{j} \binom{n}{k-j} = \binom{m+n}{k}.
\]
\textit{Hint:} Apply the binomial theorem to $(1+x)^n(1+x)^m$. % [reconstructed]
\item Prove that
\[
\binom{n}{0} - \binom{n}{1} + \binom{n}{2} - \cdots \pm \binom{n}{n} = 0.
\]
\end{letters}

\pnum{5.}
\begin{letters}
\item Prove by induction on $n$ that
\[ 1 + r + r^2 + \cdots + r^n = \frac{1 - r^{n+1}}{1 - r} \]
if $r \neq 1$ (if $r = 1$, evaluating the sum certainly presents no problem).
\item Derive this result by setting $S = 1 + r + \cdots + r^n$, multiplying
  this equation by $r$, and solving the two equations for $S$.
\end{letters}

\pnum{6.} The formula for $1^2 + \cdots + n^2$ may be derived as follows. We
begin with the formula
\[ (k+1)^3 - k^3 = 3k^2 + 3k + 1. \]
Writing this formula for $k = 1, \dots, n$ and adding, we obtain
\begin{align*}
2^3 - 1^3 &= 3 \cdot 1^2 + 3 \cdot 1 + 1\\
3^3 - 2^3 &= 3 \cdot 2^2 + 3 \cdot 2 + 1\\
&\vdots\\
(n+1)^3 - n^3 &= 3 \cdot n^2 + 3 \cdot n + 1\\
(n+1)^3 - 1 &= 3[1^2 + \cdots + n^2] + 3[1 + \cdots + n] + n.
\end{align*}
Thus we can find $\sum_{k=1}^{n} k^2$ if we already know $\sum_{k=1}^{n} k$
(which could have been found in a similar way). Use this method to find
\begin{romans}
\item $1^3 + \cdots + n^3$.
\item $1^4 + \cdots + n^4$.
\item $\dfrac{1}{1\cdot 2} + \dfrac{1}{2\cdot 3} + \cdots + \dfrac{1}{n(n+1)}$.
\item $\dfrac{3}{1^2 \cdot 2^2} + \dfrac{5}{2^2 \cdot 3^2} + \cdots +
  \dfrac{2n+1}{n^2(n+1)^2}$.
\end{romans}

\pnum{*7.} Use the method of Problem 6 to show that $\sum_{i=1}^{n} k^p$ can
always be written in the form
\[ \frac{n^{p+1}}{p+1} + A n^p + B n^{p-1} + C n^{p-2} + \cdots . \]
(The first 10 such expressions are
\[
\setlength{\arraycolsep}{3pt}
\begin{array}{ll}
\displaystyle\sum_{k=1}^{n} k = \frac{1}{2}n^2 + \frac{1}{2}n
& \displaystyle\sum_{k=1}^{n} k^6 = \frac{1}{7}n^7 + \frac{1}{2}n^6
  + \frac{1}{2}n^5 - \frac{1}{6}n^3 + \frac{1}{42}n\\[1em]
\displaystyle\sum_{k=1}^{n} k^2 = \frac{1}{3}n^3 + \frac{1}{2}n^2 + \frac{1}{6}n
& \displaystyle\sum_{k=1}^{n} k^7 = \frac{1}{8}n^8 + \frac{1}{2}n^7
  + \frac{7}{12}n^6 - \frac{7}{24}n^4 + \frac{1}{12}n^2\\[1em]
\displaystyle\sum_{k=1}^{n} k^3 = \frac{1}{4}n^4 + \frac{1}{2}n^3 + \frac{1}{4}n^2
& \displaystyle\sum_{k=1}^{n} k^8 = \frac{1}{9}n^9 + \frac{1}{2}n^8
  + \frac{2}{3}n^7 - \frac{7}{15}n^5 + \frac{2}{9}n^3 - \frac{1}{30}n\\[1em]
\displaystyle\sum_{k=1}^{n} k^4 = \frac{1}{5}n^5 + \frac{1}{2}n^4 + \frac{1}{3}n^3
  - \frac{1}{30}n
& \displaystyle\sum_{k=1}^{n} k^9 = \frac{1}{10}n^{10} + \frac{1}{2}n^9
  + \frac{3}{4}n^8 - \frac{7}{10}n^6 + \frac{1}{2}n^4 - \frac{3}{20}n^2\\[1em]
\displaystyle\sum_{k=1}^{n} k^5 = \frac{1}{6}n^6 + \frac{1}{2}n^5 + \frac{5}{12}n^4
  - \frac{1}{12}n^2
& \displaystyle\sum_{k=1}^{n} k^{10} = \frac{1}{11}n^{11} + \frac{1}{2}n^{10}
  + \frac{5}{6}n^9 - n^7 + n^5 - \frac{1}{2}n^3 + \frac{5}{66}n
\end{array}
\]
Notice that the coefficients in the second column are always $\tfrac{1}{2}$,
and that after the third column the powers of $n$ with nonzero coefficients
decrease by 2 until $n^2$ or $n$ is reached. The coefficients in all but the
first two columns seem to be rather haphazard, but there actually is some
sort of pattern; finding it may be regarded as a super-perspicacity test. See
Problem 27-17 for the complete story.)

\pnum{8.} Prove that every natural number is either even or odd.

\pnum{9.} Prove that if a set $A$ of natural numbers contains $n_0$ and
contains $k+1$ whenever it contains $k$, then $A$ contains all natural
numbers $\geq n_0$.

\pnum{10.} Prove the principle of mathematical induction from the well-ordering
principle.

\pnum{11.} Prove the principle of complete induction from the ordinary
principle of induction. \textit{Hint:} If $A$ contains 1 and $A$ contains
$n+1$ whenever it contains $1, \dots, n$, consider the set $B$ of all $k$
such that $1, \dots, k$ are all in $A$.

\pnum{12.}
\begin{letters}
\item If $a$ is rational and $b$ is irrational, is $a+b$ necessarily
  irrational? What if $a$ and $b$ are both irrational?
\item If $a$ is rational and $b$ is irrational, is $ab$ necessarily irrational?
  (Careful!)
\item Is there a number $a$ such that $a^2$ is irrational, but $a^4$ is
  rational?
\item Are there two irrational numbers whose sum and product are both
  rational?
\end{letters}

\pnum{13.}
\begin{letters}
\item Prove that $\sqrt{3}$, $\sqrt{5}$, and $\sqrt{6}$ are irrational.
  \textit{Hint:} To treat $\sqrt{3}$, for example, use the fact that every
  integer is of the form $3n$ or $3n+1$ or $3n+2$. Why doesn't this proof
  work for $\sqrt{4}$?
\item Prove that $\sqrt[3]{2}$ and $\sqrt[3]{3}$ are irrational.
\end{letters}

\pnum{14.} Prove that
\begin{letters}
\item $\sqrt{2} + \sqrt{6}$ is irrational.
\item $\sqrt{2} + \sqrt{3}$ is irrational.
\end{letters}

\pnum{15.}
\begin{letters}
\item Prove that if $x = p + \sqrt{q}$ where $p$ and $q$ are rational, and
  $m$ is a natural number, then $x^m = a + b\sqrt{q}$ for some rational $a$
  and $b$.
\item Prove also that $(p - \sqrt{q})^m = a - b\sqrt{q}$.
\end{letters}

\pnum{16.}
\begin{letters}
\item Prove that if $m$ and $n$ are natural numbers and $m^2/n^2 < 2$, then
  $(m+2n)^2/(m+n)^2 > 2$; show, moreover, that
\[ \frac{(m+2n)^2}{(m+n)^2} - 2 < 2 - \frac{m^2}{n^2}. \]
\item Prove the same results with all inequality signs reversed.
\item Prove that if $m/n < \sqrt{2}$, then there is another rational number
  $m'/n'$ with $m/n < m'/n' < \sqrt{2}$.
\end{letters}

\pnum{*17.} It seems likely that $\sqrt{n}$ is irrational whenever the natural
number $n$ is not the square of another natural number. Although the method
of Problem 13 may actually be used to treat any particular case, it is not
clear in advance that it will always work, and a proof for the general case
requires some extra information. A natural number $p$ is called a prime
number if it is impossible to write $p = ab$ for natural numbers $a$ and $b$
unless one of these is $p$, and the other 1; for convenience we also agree
that 1 is not a prime number. The first few prime numbers are
$2, 3, 5, 7, 11, 13, 17, 19$. If $n > 1$ is not a prime, then $n = ab$, with
$a$ and $b$ both $< n$; if either $a$ or $b$ is not a prime it can be factored
similarly; continuing in this way proves that we can write $n$ as a product
of primes. For example, $28 = 4 \cdot 7 = 2 \cdot 2 \cdot 7$.

\begin{letters}
\item Turn this argument into a rigorous proof by complete induction. (To be
  sure, any reasonable mathematician would accept the informal argument, but
  this is partly because it would be obvious to her how to state it
  rigorously.)

A fundamental theorem about integers, which we will not prove here, states
that this factorization is unique, except for the order of the factors. Thus,
for example, 28 can never be written as a product of primes one of which is
3, nor can it be written in a way that involves 2 only once (now you should
appreciate why 1 is not allowed as a prime).
\item Using this fact, prove that $\sqrt{n}$ is irrational unless $n = m^2$
  for some natural number $m$.
\item Prove more generally that $\sqrt[k]{n}$ is irrational unless $n = m^k$.
\item No discussion of prime numbers should fail to allude to Euclid's
  beautiful proof that there are infinitely many of them. Prove that there
  cannot be only finitely many prime numbers $p_1, p_2, p_3, \dots, p_n$ by
  considering $p_1 \cdot p_2 \cdots p_n + 1$.
\end{letters}

\pnum{*18.}
\begin{letters}
\item Prove that if $x$ satisfies
\[ x^n + a_{n-1} x^{n-1} + \cdots + a_0 = 0, \]
for some integers $a_{n-1}, \dots, a_0$, then $x$ is irrational unless $x$ is
an integer. (Why is this a generalization of Problem 17?)
\item Prove that $\sqrt{6} - \sqrt{2} - \sqrt{3}$ is irrational.
\item Prove that $\sqrt{2} + \sqrt[3]{2}$ is irrational. \textit{Hint:} Start
  by working out the first 6 powers of this number.
\end{letters}

\pnum{19.} Prove Bernoulli's inequality: If $h > -1$, then
\[ (1+h)^n \geq 1 + nh \]
for any natural number $n$. Why is this trivial if $h > 0$?

\pnum{20.} The Fibonacci sequence $a_1, a_2, a_3, \dots$ is defined as
follows:
\begin{align*}
a_1 &= 1,\\
a_2 &= 1,\\
a_n &= a_{n-1} + a_{n-2} \qquad \text{for } n \geq 3.
\end{align*}
This sequence, which begins $1, 1, 2, 3, 5, 8, \dots$, was discovered by
Fibonacci (circa 1175--1250), in connection with a problem about rabbits.
Fibonacci assumed that an initial pair of rabbits gave birth to one new pair
of rabbits per month, and that after two months each new pair behaved
similarly. The number $a_n$ of pairs born in the $n$th month is
$a_{n-1} + a_{n-2}$, because a pair of rabbits is born for each pair born the
previous month, and moreover each pair born two months ago now gives birth to
another pair. The number of interesting results about this sequence is truly
amazing---there is even a Fibonacci Association which publishes a journal,
\textit{The Fibonacci Quarterly}. Prove that
\[
a_n = \frac{\left(\dfrac{1+\sqrt{5}}{2}\right)^{n}
     - \left(\dfrac{1-\sqrt{5}}{2}\right)^{n}}{\sqrt{5}}.
\]
% [reconstructed — the closed form is the Binet formula; verify against source]
One way of deriving this astonishing formula is presented in Problem 24-16.

\pnum{21.} The Schwarz inequality (Problem 1-19) actually has a more general
form:
\[
\sum_{i=1}^{n} x_i y_i
\ \leq\ \sqrt{\sum_{i=1}^{n} x_i^{\,2}}\ \sqrt{\sum_{i=1}^{n} y_i^{\,2}}.
\]
Give three proofs of this, analogous to the three proofs in Problem 1-19.

\pnum{22.} The result in Problem 1-7 has an important generalization: If
$a_1, \dots, a_n \geq 0$, then the ``arithmetic mean''
\[ A_n = \frac{a_1 + \cdots + a_n}{n} \]
and ``geometric mean''
\[ G_n = \sqrt[n]{a_1 a_2 \cdots a_n} \]
satisfy
\[ G_n \leq A_n. \]

\begin{letters}
\item Suppose that $a_1 < A_n$. Then some $a_i$ satisfies $a_i > A_n$; for
  convenience, say $a_2 > A_n$. Let $\bar{a}_1 = A_n$ and let
  $\bar{a}_2 = a_1 + a_2 - \bar{a}_1$. Show that
\[ \bar{a}_1 \bar{a}_2 \geq a_1 a_2. \]
Why does repeating this process enough times eventually prove that
$G_n \leq A_n$? (This is another place where it is a good exercise to provide
a formal proof by induction, as well as an informal reason.) When does
equality hold in the formula $G_n \leq A_n$?

The reasoning in this proof is related to another interesting proof.
\item Using the fact that $G_n \leq A_n$ when $n = 2$, prove, by induction on
  $k$, that $G_n \leq A_n$ for $n = 2^k$.
\item For a general $n$, let $2^m > n$. Apply part (b) to the $2^m$ numbers
\[ a_1, \dots, a_n, \underbrace{A_n, \dots, A_n}_{2^m - n \text{ times}} \]
to prove that $G_n \leq A_n$.
\end{letters}

\pnum{23.} The following is a recursive definition of $a^n$:
\begin{align*}
a^1 &= a,\\
a^{n+1} &= a^n \cdot a.
\end{align*}
Prove, by induction, that
\[ a^{n+m} = a^n \cdot a^m, \qquad (a^n)^m = a^{nm}. \]
(Don't try to be fancy: use either induction on $n$ or induction on $m$, not
both at once.)

\pnum{24.} Suppose we know properties P1 and P4 for the natural numbers, but
that multiplication has never been mentioned. Then the following can be used
as a recursive definition of multiplication:
\begin{align*}
1 \cdot b &= b,\\
(a+1) \cdot b &= a \cdot b + b.
\end{align*}
Prove the following (in the order suggested!):
\[
a \cdot (b+c) = a \cdot b + a \cdot c \quad\text{(use induction on }a\text{)},
\]
\[
a \cdot 1 = a,
\]
\[
a \cdot b = b \cdot a \quad\text{(you just finished proving the case }b=1\text{)}.
\]

\pnum{25.} In this chapter we began with the natural numbers and gradually
built up to the real numbers. A completely rigorous discussion of this
process requires a little book in itself (see Part V). No one has ever
figured out how to get to the real numbers without going through this
process, but if we do accept the real numbers as given, then the natural
numbers can be defined as the real numbers of the form $1, 1+1, 1+1+1$, etc.
The whole point of this problem is to show that there is a rigorous
mathematical way of saying ``etc.''

\begin{letters}
\item A set $A$ of real numbers is called \emph{inductive} if
\begin{enumerate}[label=(\arabic*),leftmargin=2.6em,itemsep=0pt,topsep=2pt]
\item $1$ is in $A$,
\item $k+1$ is in $A$ whenever $k$ is in $A$.
\end{enumerate}
Prove that
\begin{romans}
\item $\mathbb{R}$ is inductive.
\item The set of positive real numbers is inductive.
\item The set of positive real numbers unequal to $\tfrac{1}{2}$ is inductive.
\item The set of positive real numbers unequal to 5 is not inductive.
\item If $A$ and $B$ are inductive, then the set $C$ of real numbers which
  are in both $A$ and $B$ is also inductive.
\end{romans}
\item A real number $n$ will be called a \emph{natural number} if $n$ is in
  every inductive set.
\begin{romans}
\item Prove that 1 is a natural number.
\item Prove that $k+1$ is a natural number if $k$ is a natural number.
\end{romans}
\end{letters}

\pnum{26.} There is a puzzle consisting of three spindles, with $n$ concentric
rings of decreasing diameter stacked on the first (Figure 1). A ring at the
top of a stack may be moved from one spindle to another spindle, provided
that it is not placed on top of a smaller ring. For example, if the smallest
ring is moved to spindle 2 and the next-smallest ring is moved to spindle 3,
then the smallest ring may be moved to spindle 3 also, on top of the
next-smallest. Prove that the entire stack of $n$ rings can be moved onto
spindle 3 in $2^n - 1$ moves, and that this cannot be done in fewer than
$2^n - 1$ moves.

% Figure 1: the three-spindles puzzle. Semantic TikZ (ellipses, cylinders,
% TeX labels) — reference scan: figures/scans/calcu5-fig1.png.
\begin{center}
% Figure 1 — three-spindle puzzle (Problem 2-26).
% Semantic reconstruction of figures/scans/calcu5-fig1.png.
% Oblique board with mild perspective (front a little wider than the back,
% left edge more diagonal, right edge steeper). Four disks + spindle;
% two empty pegs. Labels and caption are TeX.

\begin{tikzpicture}[reconstruction, line cap=round, line join=round,
  line width=0.9pt, x=1mm, y=1mm]

  % Board corners — match the scan's nearly-horizontal front/back.
  \coordinate (FL) at (0.0, 2.0);
  \coordinate (FR) at (106.0, 6.0);
  \coordinate (BL) at (18.0, 50.0);
  \coordinate (BR) at (116.0, 52.5);

  % Peg 2 sits on the back edge (scan: the edge is interrupted, not drawn through).
  \def\pxtwo{70.0}
  \def\pytwo{51.3}
  \def\prad{5.5}
  \def\pry{2.1}

  \draw (FL) -- (FR) -- (BR);
  \draw (FL) -- (BL);
  \draw (BL) -- ({\pxtwo-\prad}, {\pytwo});
  \draw ({\pxtwo+\prad}, {\pytwo}) -- (BR);

  % Empty peg: opaque white, two ellipses, two generators.
  \newcommand{\hanoipeg}[5]{% #1,#2 base centre; #3 rx; #4 ry; #5 height
    \fill[white]
      ({#1-#3}, {#2}) -- ({#1-#3}, {#2+#5})
      arc[start angle=180, end angle=0, x radius=#3, y radius=#4]
      -- ({#1+#3}, {#2})
      arc[start angle=0, end angle=-180, x radius=#3, y radius=#4]
      -- cycle;
    \fill[white] ({#1}, {#2}) ellipse [x radius=#3, y radius=#4];
    \fill[white] ({#1}, {#2+#5}) ellipse [x radius=#3, y radius=#4];
    \draw ({#1}, {#2}) ellipse [x radius=#3, y radius=#4];
    \draw ({#1-#3}, {#2}) -- ({#1-#3}, {#2+#5});
    \draw ({#1+#3}, {#2}) -- ({#1+#3}, {#2+#5});
    \draw ({#1}, {#2+#5}) ellipse [x radius=#3, y radius=#4];
  }

  % Disk: shaded cylindrical wall (light from the left, as in the scan),
  % white top, full bottom ellipse.
  \newcommand{\hanoidisk}[5]{% #1,#2 top centre; #3 rx; #4 ry; #5 thickness
    \shade[left color=black!18, right color=black!38]
      ({#1-#3}, {#2-#5}) -- ({#1-#3}, {#2})
      arc[start angle=180, end angle=360, x radius=#3, y radius=#4]
      -- ({#1+#3}, {#2-#5})
      arc[start angle=360, end angle=180, x radius=#3, y radius=#4]
      -- cycle;
    \draw ({#1-#3}, {#2}) -- ({#1-#3}, {#2-#5});
    \draw ({#1+#3}, {#2}) -- ({#1+#3}, {#2-#5});
    \draw ({#1}, {#2-#5}) ellipse [x radius=#3, y radius=#4];
    \fill[white] ({#1}, {#2}) ellipse [x radius=#3, y radius=#4];
    \draw ({#1}, {#2}) ellipse [x radius=#3, y radius=#4];
  }

  % Peg 2 — on the rear edge, tall. Peg 3 is further right and forward;
  % the scan leaves a clear gap between them.
  \hanoipeg{\pxtwo}{\pytwo}{\prad}{\pry}{31}
  \node[font=\large] at (\pxtwo, 41.5) {2};

  \hanoipeg{98.0}{18.5}{\prad}{\pry}{28}
  \node[font=\large] at (98.0, 9.5) {3};

  % Stack, bottom disk first. Four rings + a short spindle (same radius as the pegs).
  \hanoidisk{37}{24.0}{21.5}{8.0}{4.4}
  \hanoidisk{37}{28.0}{17.4}{6.5}{4.0}
  \hanoidisk{37}{31.6}{13.6}{5.05}{3.6}
  \hanoidisk{37}{34.8}{10.0}{3.7}{3.2}
  \hanoipeg{37}{34.8}{5.5}{2.1}{6.8}

  \node[anchor=west, font=\normalsize, inner sep=0pt]
    at (0, -3.0) {\textls[260]{FIGURE 1}};
\end{tikzpicture}

\end{center}

\pnum{*27.} University B. once boasted 17 tenured professors of mathematics.
Tradition prescribed that at their weekly luncheon meeting, faithfully
attended by all 17, any members who had discovered an error in their
published work should make an announcement of this fact, and promptly resign.
Such an announcement had never actually been made, because no professor was
aware of any errors in her or his work. This is not to say that no errors
existed, however. In fact, over the years, in the work of every member of the
department at least one error had been found, by some other member of the
department. This error had been mentioned to all other members of the
department, but the actual author of the error had been kept ignorant of the
fact, to forestall any resignations.

One fateful year, the department was augmented by a visitor from another
university, one Prof.~X, who had come with hopes of being offered a permanent
position at the end of the academic year. Naturally, he was apprised, by
various members of the department, of the published errors which had been
discovered. When the hoped-for appointment failed to materialize, Prof.~X
obtained his revenge at the last luncheon of the year. ``I have enjoyed my
visit here very much,'' he said, ``but I feel that there is one thing that I
have to tell you. At least one of you has published an incorrect result,
which has been discovered by others in the department.'' What happened the
next year?

\pnum{**28.} After figuring out, or looking up, the answer to Problem 27,
consider the following: Each member of the department already knew what
Prof.~X asserted, so how could his saying it change anything?
